\documentclass[10pt,a4paper]{amsart}
\usepackage[margin=19mm]{geometry}
\usepackage{amsmath,amssymb,mathtools}
\usepackage[scr=rsfs,cal=euler]{mathalfa}
\usepackage{xfrac}
\usepackage[unicode,hidelinks,bookmarks=false]{hyperref}
\newcommand{\ie}{i.e.\ }
\newcommand{\cf}{cf.\ }
\newcommand{\resp}{respectively\ }
\newcommand{\Subsection}{\subsection}
\providecommand{\nobreakdash}{\nobreak\hbox{-}}
\setlength{\emergencystretch}{2em}
\multlinegap0pt
\usepackage{bm}
\usepackage{bbm}
\usepackage{dsfont}
\usepackage{xspace}
\usepackage{cancel}
\usepackage{mathtools}
\usepackage{amsthm}
\makeatletter
\@ifundefined{thm}{\newtheorem{thm}{Theorem}}{}
\makeatother
\title[Harder-Narasimhan Filtration on Moment Map for Quiver Repres]{Harder-Narasimhan Filtration on Moment Map for Quiver Representations}
\author{Ching Yan Timothy Yau}
\date{Source: arXiv:2511.02449v2 (CC BY 4.0)}
\begin{document}
\maketitle

\noindent\emph{Source and licence.} Harder-Narasimhan Filtration on Moment Map for Quiver Representations. Version arXiv:2511.02449v2 (CC BY 4.0), distributed under the Creative Commons Attribution 4.0 International licence, as recorded in the arXiv licence metadata for this exact version and shown on the abstract page https://arxiv.org/abs/2511.02449v2. Full rights notice: licenses/arxiv-2511.02449-CC-BY-4.0.txt.\par\smallskip

\noindent\textit{Editorial scope.} Counted unit: the Thm whose source label is \texttt{None} in the article (source0.tex lines 330--331, source label \texttt{None}), delivered together with the article's own complete original proof of it (source0.tex lines 332--356, one \texttt{proof} environment of 25 lines). No mathematics is written, repaired or re-derived: statement and proof are the article's own text line for line. Required context retained uncounted: none. Every definition, display and label of those lines is retained unchanged. Omitted from this package and not redistributed: 1--329, 357--588. Nothing inside a retained excerpt is omitted or altered: every retained body is delivered exactly as the article's lines are written, so no omission receipt of any kind accompanies this package. Citations: the retained excerpts carry no citation command at all, so no bibliography entry of the article is transcribed and no reference is redistributed. Reference adaptation: none was needed and none was made; every reference inside the delivered unit resolves inside it, so no unresolved reference or citation is produced. Printed slips kept literal: no printed word, symbol, hypothesis or number of the counted statement or of its proof was altered. Layout and class: the article's own class is \texttt{amsart} with packages including amscd, amsxtra, amssymb, mathrsfs, bbm, url, amsmath, geometry, tikz-cd; the wrapper is the lane's pinned \texttt{amsart} editorial wrapper at 10pt on A4, which restates the theorem-like environments the retained text needs and adds the article's own macro definitions that the retained text needs (none), with extra wrapper packages (bm, bbm, dsfont, xspace, cancel, mathtools). The delivered numbering is the unit's own. Assets: no retained span contains a figure environment, a graphics-inclusion command, a PDF-page inclusion, a table or a TikZ picture; no image file is copied into this package, the package declares no asset dependency and no image file is redistributed with it. Licence: version arXiv:2511.02449v2 of this article is distributed under the Creative Commons Attribution 4.0 International licence, as recorded in the arXiv licence metadata for this exact version and shown on the abstract page https://arxiv.org/abs/2511.02449v2. A full rights notice is delivered at licenses/arxiv-2511.02449-CC-BY-4.0.txt. Characters: every retained excerpt is pure ASCII, so the delivered engine, XeLaTeX with its default Latin Modern text font, reports no missing character.
\begin{thm} If $\epsilon_{F^*}\geq n$, then $s_0=1$ and $s_l=0$ for all $1\leq l\leq n-1$ as in definition 4.1.
\end{thm}

\begin{proof}
We proceed by induction on $n$. The case $n=0$ is evident.
\\~\\ Suppose all smaller cases hold, then by inductive hypothesis it suffices to prove $s_{n-1}\leq1$, if $(g_i)_{i\in I}\in S_{n-1}$, then by inductive hypothesis there exists $i,j\in I$ such that in the notation of Lemma 5.1, $c(g_i,g_j)\geq \dim Hom(V_i,V_j)_{F^*}$. 
We first notice that the element whose $i$ component is $\lambda \operatorname{id}$ for every $i\in I$ belongs to $S_0$ and hence all of $S_1,...,S_{n-1}$, hence $s_0,...,s_{n-1}\geq 1$. Therefore, it suffices to show that $\beta:i\to j$ is an edge, then for $A\in End(V)_{D(F^*)},B\in End(W)_{D(F^*)}$ where $(A,B)\neq (\lambda \operatorname{id},\lambda\operatorname{id})$, then we must have
$$c(A,B)\leq \dim Hom(V,W)_{F^*}-\epsilon_{F^*}\hspace{20pt}(1)$$
let $A=l_Au_A$ and $B=l_Bu_B$ be its Levi decomposition in its respective parabolic subalgebra $End(V)_{F^*},End(W)_{F^*}$. We notice that if we conjugate $A$ by some $p=lv\in End(V)_{F^*}$, then the levi decomposition of $pAp^{-1}$ will be $ll_Al^{-1}$. Together with the fact that $c(pAp^{-1},B)=c(A,B)$ we can WLOG assume that $A$, and similarly $B$ are both lower triangular matrices.
\\~\\ Now we let $D_A$ and $D_B$ be the set of diagonal entries of $A,B$ respectively. We divide into two cases:
\\~\\\underline{Case I:} $|D_{A}\cup D_B|\geq 2$
\\ Pick positive integers $a_1>a_2>...>a_{\dim V}$. Define, for all $t\in\mathbb C$, matrices $M_t=\operatorname{diag}(t^{a_1},...,t^{a_n})$. Consider the mapping $\lambda:\mathbb C^*\to End(V)_{*}$ 
$$t\mapsto M_tAM_t^{-1}$$
We notice that 
$$(M_tAM_t^{-1})_{ij}=(t^{a_i-a_j}A)_{ij}$$
By the assumption that $A$ is upper triangular, we notice that the matrix $d(A)$ consisting only of the diagonal entries of $A$ lies in $\lambda(\mathbb C^*)$, meanwhile, we have that $c(M_tAM_t^{-1},B)=c(A,B)$.
Therefore, by Lemma 5.1, we have 
$$c(d(A),B)\geq c(A,B)$$
Hence, we may WLOG assume $A$ and similarly $B$ are diagonal matrices. Now  either
\begin{enumerate}
    \item[(i)] For each $1\leq j\leq \epsilon_{F^*}$ there exists $1\leq i\leq \dim V$ such that $a_{i,i}\neq b_{j,j}$, or
    \item[(ii)] for each $1\leq i\leq \epsilon_{F^*}$ there exists some $1\leq j\leq \dim W$ such that $a_{\dim V-i,\dim V-i}\neq b_{j,j}$
\end{enumerate} 
In the former case, the map sending $v_i$ to $w_j$ is not in $r(A,B)$, and in the latter case, the map sending $v_{\dim V-i}$ to $w_j$ is not in $r(A,B)$. However, both of these are in $Hom(V,W)_{F^*}$ by the definition of $\epsilon{F^*}$. Hence, the codimension of $r(A,B)$ is at least $d$.
\\~\\ \underline{Case II}: $|D_A\cup D_B|=1$
\\ This is equivalent to all the diagonal entries of $A,B$ being the same. Therefore, we may write $A=\lambda I+A',B=\lambda I+B'$. Notice that $\lambda I$ commutes with all maps, so $$c(A,B)=c(A',B')$$ By our assumption, at least one of $A',B'$ must be nonzero. Suppose $A'\neq 0$, then by the $\mathbb C^*$ action argument in Case I, we may assume $B'=0$. Then we notice that $fA=0$ if the image of $A$, which is nonzero by assumption, lies in the kernel of $f$. This again shows that the codimension of $r(A,B)$ is at least $\epsilon_{F^*}$. 
\\~\\ Similarly, if $B'\neq 0$, then we can assume $A'=0$, then $B'f=0$ implies that $f$ must map into the kernel of $B'$, which is not $W$ by assumption, and hence the codimension is at least $\epsilon_{F^*}$ by the same argument.
\end{proof}

\end{document}
